%% TRANSFORMATIONS TEMPORELLES %%
%
\section{Sommation}

\itemstyle{ Op\'erateur~: } {\sf Som}
\index{Som}

%%%%%OBJET%%%%%
\itemstyle{ Objet~: }
Sommation de $N$ \'elements cons\'ecutifs dans le temps.\newline


%%%%%ARGUMENTS%%%%%
\bigskip
%%%%%ENTREES%%%%
\begin{tabarg}
\multicolumn{3}{l}{\titretab{ Entr\'ees}}\\
\hline
$X$ & Vecteur, {\tt float} & Param\`etre d'entr\'ee\\
$N$ & Coeff., {\tt int} & Nombre de valeurs de la sommation \\
$Decal$ & Coeff., {\tt int} & D\'ecalage (\textit{nombre de mesures})\\
\end{tabarg}
\medskip

%%%%%SORTIES%%%%
\begin{tabarg}
\multicolumn{3}{l}{\titretab{ Sorties}}\\
\hline
$Xs$ & Vecteur, {\tt float} & Param\`etre r\'esultat.\\
\end{tabarg}


%%%%%FORMULE%%%%
\itemstyle{ Formule }

%%Som
\[
Xs[i] = \sum_{j= N\times i + Decal}^{N\times (i+1) -1 + Decal} X[j] 
 \qquad i=0, Nbre\_echant(X)/N -1 
\]



%%%%%COMMENTAIRE%%%%
\itemstyle{ Commentaire~:}

Le nombre d'\'echantillons du vecteur r\'esultat est \'egal au nombre
d'\'echantillons du vecteur d'origine, divis\'e par  $N$, arrondi \`a
l'entier sup\'erieur.
Si le nombre d'\'echantillons de X n'est pas un multiple de $N$, la
derni\`ere valeur de $Xs$ est identique � la pr�c�dente.\\
La fr\'equence d'\'echantillonnage du param\`etre r\'esultat est le rapport de
la fr\'equence du param\`etre d'entr\'ee par $N$.\\ 
Cas particulier :\\
Si le nombre d'\'echantillons de X est inf�rieur � $N$, on renvoie un NaN (modification du source en date du 26/10/2007)

%%%%%AUTEUR%%%%%
\itemstyle{ Auteur~:}
CNRM/GMEI/TRAMM

%%%%%REFERENCES%%%%%
\itemstyle{ R\'ef\'erences~:}
